% Generated draft: requires mathematical and source review.
\documentclass[11pt]{article}
\usepackage[margin=25mm]{geometry}
\usepackage{fontspec}
\setmainfont{Latin Modern Roman}
\usepackage{amsmath,amssymb,mathtools,unicode-math}
\setmathfont{Latin Modern Math}
\usepackage{enumitem,needspace,xurl,hyperref}
\setlength{\parindent}{0pt}
\setlength{\parskip}{5pt}
\setlength{\emergencystretch}{3em}
\newcommand{\N}{\mathbb{N}}
\newcommand{\Z}{\mathbb{Z}}
\newcommand{\Q}{\mathbb{Q}}
\newcommand{\R}{\mathbb{R}}
\newcommand{\C}{\mathbb{C}}
\newcommand{\F}{\mathbb{F}}
\newcommand{\A}{\mathbb{A}}
\newcommand{\PP}{\mathbb{P}}
\newcommand{\E}{\mathbb{E}}
\newcommand{\eps}{\varepsilon}
\newcommand{\dd}{\,\mathrm{d}}
\begin{document}
\section*{UnsolvedMath ID: 30005592}
\textbf{Source problem code:} catalogue entry
\subsection{Definitions and mathematical statement}
A four-dimensional steady gradient Ricci soliton is a triple \((M^4,g,f)\), where \(g\) is a smooth Riemannian metric and \(f\colon M\to\R\) is smooth, satisfying
\[
\operatorname{Ric}_g+\nabla_g^2 f=0.
\]
This fixes the sign convention for the potential; replacing \(f\) by \(-f\) gives the alternative convention. The curvature operator is the self-adjoint map \(\mathcal R_g\colon\Lambda^2T_xM\to\Lambda^2T_xM\) induced by the Riemann curvature tensor. Nonnegative curvature operator means
\[
\langle\mathcal R_g\omega,\omega\rangle_g\geq0
\qquad(x\in M,\ \omega\in\Lambda^2T_xM).
\]

The source also requires the soliton to be \emph{noncollapsed}, but does not specify its definition or scale range. One common metric interpretation is that there exists \(\kappa>0\) such that, for every \(x\in M\) and \(r>0\),
\[
\sup_{B_g(x,r)}|\operatorname{Rm}_g|\leq r^{-2}
\quad\Longrightarrow\quad
\operatorname{Vol}_g(B_g(x,r))\geq\kappa r^4.
\]
This formula is a possible convention for review, not an additional hypothesis established by the supplied record; an interpretation using the associated Ricci flow and parabolic neighborhoods would require its own specification.

Let \(B^4\) denote the four-dimensional Bryant steady gradient Ricci soliton. Let \(\mathcal F_{[5]}\) provisionally denote the particular family described in the original excerpt as the \(\Z_2\times O(3)\)-symmetric solitons constructed in reference [5]. Here \(\Z_2\) is the cyclic group of order two, and \(O(3)=\{A\in\R^{3\times3}:A^{\mathsf T}A=I\}\). Symmetry refers to an isometric group action, but the intended action, its effectiveness, its action on the potential, and the defining construction and parameters of \(\mathcal F_{[5]}\) are not supplied. In particular, \(\mathcal F_{[5]}\) cannot be defined merely as all solitons having this symmetry.

\textbf{Source conjecture, with unresolved scope.} Every noncollapsed four-dimensional steady gradient Ricci soliton with nonnegative curvature operator belongs to the Bryant model or to \(\mathcal F_{[5]}\). The source does not specify whether membership is understood up to isometry, constant metric rescaling, or another equivalence of soliton triples. It also does not explicitly impose completeness, noncompactness, nonflatness, simple connectivity, or exclusion of product solitons and quotients. These omissions prevent an unqualified precise classification statement; none of these hypotheses is inserted here.

The second assertion is that every member of \(\mathcal F_{[5]}\) has a ray as its blow-down. A ray means the metric space \(([0,\infty),d_{\mathrm{Euc}})\). For a metric interpretation, a pointed blow-down at \(p\in M\) would be a pointed Gromov--Hausdorff limit
\[
(M,\lambda_j g,p)\longrightarrow(X,d_X,p_X),
\qquad \lambda_j>0,\quad\lambda_j\longrightarrow0.
\]
The intended conclusion would identify the limit with
\[
(X,d_X,p_X)\cong([0,\infty),d_{\mathrm{Euc}},0).
\]
However, the supplied record does not specify the convergence topology, basepoint convention, existence of limits, or whether the assertion concerns one subsequential limit, every subsequential limit, or convergence along every sequence of scales. Thus this display records a candidate interpretation of the claimed ray limit, not a resolved quantification of the source assertion.

\textbf{Additional questions in the original excerpt.} The excerpt also asks about \(O(2)\)-symmetry for three-dimensional ancient Ricci flows and for higher-dimensional steady gradient solitons with positive curvature operator, and about \(O(3)\)-symmetry for all noncollapsed four-dimensional steady gradient solitons. An ancient Ricci flow satisfies \(\partial_tg(t)=-2\operatorname{Ric}_{g(t)}\) on a time interval extending to \(-\infty\); positive curvature operator means strict positivity on every nonzero two-form. The excerpt supplies no further hypotheses or precise meaning of these symmetry assertions. The four-dimensional symmetry question does not explicitly include the nonnegative-curvature-operator assumption of the main conjecture. These questions are retained as contextual questions, not silently incorporated as established statements or additional hypotheses.
\subsection{Short English statement}
The supplied conjecture seeks to classify noncollapsed four-dimensional steady gradient Ricci solitons with nonnegative curvature operator as the Bryant soliton or members of a particular \(\Z_2\times O(3)\)-symmetric family, and asserts that members of that family have a ray as their blow-down. The family, admissible geometric hypotheses, equivalence relation, noncollapse convention, and blow-down quantifiers require clarification. This is a statement-only draft requiring human mathematical review.
\subsection{Sources}
\begin{itemize}
\item Local UnsolvedMath registry, permanent ID \texttt{30005592}, source problem code \texttt{catalogue entry}, title ``Classification of Noncollapsed Four-Dimensional Steady Solitons.'' The supplied record labels the problem open and the statement \texttt{corrected\_verified}, and states that it was reconstructed and checked against the official report. These are supplied registry assertions, not independent verification performed for this draft.
\item Supplied primary-source citation: \emph{Partial Differential Equations (2024)}, Oberwolfach Reports, DOI 10.4171/owr/2023/32, pp.\ 1789--1842, \url{https://doi.org/10.4171/owr/2023/32}. The excerpt labels the main assertion ``Conjecture 2.'' The registry gives proposed year 2023 and no proposer. The relationship between the citation's 2024 label and the 2023 provenance requires checking.
\item Supplied literature source ``Oberwolfach Report 32'': \url{https://doi.org/10.4171/owr/2023/32}.
\item Supplied literature source ``Four-dimensional partial classification'': \url{https://arxiv.org/abs/2403.19627}. No authors or full bibliographic title are supplied in the record.
\item Supplied literature source ``Law uniqueness theorem'': \url{https://arxiv.org/abs/2505.20576}. The record supplies this label without a full bibliographic citation or theorem hypotheses.
\item Literature assessment dated exactly \texttt{2026-08-22}: the supplied curation labels the problem open, describes the Bryant and constructed \(\Z_2\times O(3)\) families as known, and reports four-dimensional classification results under additional asymptotic or curvature-decay hypotheses, while describing the unrestricted classification and blow-down assertion as unproved. This is supplied dated curation, not independently verified literature assessment. Its background explicitly requests independent verification of time-sensitive status and cited claims.
\item The original excerpt attributes the symmetric family to an unnamed speaker through ``I constructed in [5].'' Reference [5] and the speaker's identity are absent. The excerpt also attributes a classification of noncollapsed translators in \(\R^4\) to Choi--Haslhofer--Hershkovits, citing [3, 2] as motivation; the corresponding bibliography entries are absent. These attributions and claims are preserved as supplied context and have not been independently verified.
\end{itemize}
\end{document}
